%% GENERATED by tools/gen_latex.py -- edit content/ch13.py instead.
%% Build:  latexmk -pdf ch13.tex     (or: pdflatex ch13.tex, twice)
\documentclass[aspectratio=169,11pt,t]{beamer}
\usepackage{itbpro}

\coursetitle{Timeseries Forecasting}
\coursesubtitle{Where a common-sense baseline beats two deep learning models in a row --- and what that teaches about architecture priors.}
\coursekicker{Chapter 13}
\coursesource{Chollet \& Watson, Deep Learning with Python 3e --- chapter 13}
\courseduration{3 hours (2 sessions)}
\coursepresenter{Prof. Bambang Riyanto Trilaksono $\cdot$ Rahman Indra Kesuma, S.Kom., M.Cs.}{Lead Instructor $\cdot$ Teaching Assistant}
\coursebrand{AI for Professional \textperiodcentered\ ITB \textperiodcentered\ Ch 13}

\title{Timeseries Forecasting}
\author{Prof. Bambang Riyanto Trilaksono $\cdot$ Rahman Indra Kesuma, S.Kom., M.Cs.}
\date{}

\begin{document}
\covertitle

\framekicker{Learning outcomes}
\begin{frame}[shrink=25]{Session Objectives}
\leadpar{By the end of this session, participants can:}
\begin{isteps}
\item Name the four common timeseries tasks and say which are supervised.
\item Build a \textbf{common-sense baseline} for a forecasting problem and evaluate it properly.
\item Explain \textbf{why a dense network and a 1D ConvNet fail} on this problem --- and what that says about architecture priors.
\item Explain what an \textbf{RNN} is, and how it differs from a feedforward network.
\item Say why \textbf{\inlinecode{SimpleRNN} is not used in practice}, and how \textbf{LSTM} solves it.
\item Apply \textbf{recurrent dropout} correctly, and say why the naive version is harmful.
\item \textbf{Stack} recurrent layers with \inlinecode{return\_sequences}, and know when a \textbf{bidirectional} layer helps and when it cannot.
\end{isteps}

\vskip 6pt
\vskip 4pt
\reslink{Book}{Chapter 13 --- full text}{https://deeplearningwithpython.io/chapters/chapter13\_timeseries-forecasting}
\reslink{Notebook}{01\_jena\_data\_and\_baseline.ipynb}{https://hendrikarisma.my.id/rs/ai-products-course/notebooks/ch13/01\_jena\_data\_and\_baseline.html}
\reslink{Notebook}{02\_dense\_and\_conv1d.ipynb}{https://hendrikarisma.my.id/rs/ai-products-course/notebooks/ch13/02\_dense\_and\_conv1d.html}
\reslink{Notebook}{03\_lstm\_and\_gru.ipynb}{https://hendrikarisma.my.id/rs/ai-products-course/notebooks/ch13/03\_lstm\_and\_gru.html}
\reslink{Notebook}{04\_dropout\_stacking\_bidirectional.ipynb}{https://hendrikarisma.my.id/rs/ai-products-course/notebooks/ch13/04\_dropout\_stacking\_bidirectional.html}
\reslink{Notebook}{All chapter 13 notebooks}{https://hendrikarisma.my.id/rs/ai-products-course/notebooks/index.html\#ch13}
\reslink{Repo}{Author's official notebook}{https://github.com/fchollet/deep-learning-with-python-notebooks/blob/master/chapter13\_timeseries-forecasting.ipynb}
\vskip 2pt
\end{frame}

\framekicker{Section 13.1}
\begin{frame}[shrink=25]{What a timeseries is, and why it is different}
{\fontsize{9.4}{12.6}\selectfont\color{fg2}Any data obtained by measurement at regular intervals: a daily stock price, a city's hourly electricity consumption, a store's weekly sales.\par}\vskip 4pt
\begin{iband}[signal]
Unlike everything so far, working with timeseries means understanding \textbf{the dynamics of a system} --- its periodic cycles, how it trends, its regular regime, and \hilite{its sudden spikes}
\end{iband}
\end{frame}

\framekicker{Section 13.1}
\begin{frame}[shrink=25]{Four things you can do with one}
\begin{center}
\includegraphics[width=0.92\linewidth,height=0.60\textheight,keepaspectratio]{figs/ch13-tasks.tex.pdf}
\end{center}
\vskip 3pt{\color{ink3}\fontsize{7.4}{9.4}\selectfont Forecasting is by far the most common, and is what this chapter covers.\par}
{\fontsize{9.4}{12.6}\selectfont\color{fg2}Note the parenthetical on anomaly detection: it is usually \textbf{unsupervised}, because you often do not know what kind of anomaly you are looking for --- so \hilite{you cannot train on examples of it}\par}\vskip 4pt
\end{frame}

\framekicker{Section 13.2}
\begin{frame}[shrink=25]{The running example}
{\fontsize{9.4}{12.6}\selectfont\color{fg2}Every code example in the chapter targets one problem: \textbf{predict the temperature 24 hours from now}, given readings from the recent past.\par}\vskip 4pt
{\fontsize{9.4}{12.6}\selectfont\color{fg2}The data is the Jena weather dataset --- fourteen meteorological quantities recorded every ten minutes over several years. It is small enough to work with and \hilite{genuinely hard enough to be instructive}\par}\vskip 4pt
\end{frame}

\framekicker{Section 13.2.1}
\begin{frame}[shrink=25]{The split has to respect time}
\begin{center}
\includegraphics[width=0.92\linewidth,height=0.60\textheight,keepaspectratio]{figs/ch13-splitsplit.tex.pdf}
\end{center}
\vskip 3pt{\color{ink3}\fontsize{7.4}{9.4}\selectfont Chapter 5's arrow-of-time pitfall, made concrete.\par}
{\fontsize{9.4}{12.6}\selectfont\color{fg2}The data must \textbf{not} be shuffled before splitting: all test data has to be later than all training data. And the normalisation statistics come from \hilite{the training portion only}\par}\vskip 4pt
\end{frame}

\framekicker{Section 13.2.1 $\cdot$ listing 13.6}
\begin{frame}[shrink=25]{Normalising fourteen quantities on different scales}
{\fontsize{9.4}{12.6}\selectfont\color{fg2}No vectorisation is needed --- the data is already numeric. But atmospheric pressure sits around \textbf{1,000 mbar} while H2OC sits around \textbf{3}, so each series is normalised independently.\par}\vskip 4pt
\icode{python}{listing 13.6 --- normalising}{listings/ch13-01.py}
\begin{iband}[signal]
Note the slice: \inlinecode{[:num\_train\_samples]}. The statistics are computed on the \textbf{first 210,225 timesteps only} --- \hilite{the same information-leak rule as chapter 4's housing example}
\end{iband}
\end{frame}

\framekicker{Section 13.2.1}
\begin{frame}[shrink=25]{Windows without copying the data}
\begin{center}
\includegraphics[width=0.92\linewidth,height=0.60\textheight,keepaspectratio]{figs/ch13-windows.tex.pdf}
\end{center}
\vskip 3pt{\color{ink3}\fontsize{7.4}{9.4}\selectfont Consecutive samples share most of their timesteps, so materialising each one would waste enormous memory.\par}
{\fontsize{9.4}{12.6}\selectfont\color{fg2}Keras has a utility for exactly this --- \inlinecode{timeseries\_dataset\_from\_array()} --- which \textbf{generates the windows on the fly}, keeping only the original array in memory.\par}\vskip 4pt
\end{frame}

\framekicker{Section 13.2.1}
\begin{frame}[shrink=25]{Understanding it on a toy example first}
{\fontsize{9.4}{12.6}\selectfont\color{fg2}The offset between \inlinecode{data} and \inlinecode{targets} is what turns windowing into forecasting, and it is easiest to see on ten integers.\par}\vskip 4pt
\icode{python}{the utility, on integers}{listings/ch13-02.py}
\ioutput{listings/ch13-03.txt}
\end{frame}

\framekicker{Section 13.2.1}
\begin{frame}[shrink=25]{\ldots{}then the real thing}
{\fontsize{9.4}{12.6}\selectfont\color{fg2}Same call, with the problem's actual parameters: five days of history sampled hourly, predicting 24 hours ahead.\par}\vskip 4pt
\icode{python}{the training dataset}{listings/ch13-04.py}
\begin{iband}[signal]
\inlinecode{start\_index} and \inlinecode{end\_index} are how the three splits are carved out of one array \textbf{without shuffling} --- \hilite{the chronological order is preserved by construction}
\end{iband}
\end{frame}

\coursesection{01}{The baseline, and two failures}{Common sense contains information a model has no access to.}

\framekicker{Section 13.2.2}
\begin{frame}[shrink=25]{Before any model: what would common sense do?}
{\fontsize{9.4}{12.6}\selectfont\color{fg2}Temperature is continuous and daily-periodic, so a sensible heuristic is: \textbf{the temperature in 24 hours will equal the temperature right now.}\par}\vskip 4pt
\icode{python}{listing 13.9 --- the common-sense baseline}{listings/ch13-05.py}
\ioutput{listings/ch13-06.txt}
\end{frame}

\framekicker{Section 13.2.2}
\begin{frame}[shrink=25]{That is the number to beat}
\begin{iband}[signal]
Assume tomorrow is like today and you are off by \textbf{two and a half degrees on average}. Not good enough to launch a forecasting service --- but \hilite{every model from here on has to beat it to justify itself}
\end{iband}
{\fontsize{9.4}{12.6}\selectfont\color{fg2}The book makes the general point too: for an unbalanced classification task with 90\% class A, \textit{always predict A} scores 90\%. \textbf{Such elementary baselines can prove surprisingly hard to beat.}\par}\vskip 4pt
\end{frame}

\framekicker{Section 13.2.3 $\cdot$ listing 13.10}
\begin{frame}[shrink=25]{Attempt 1 --- a small dense network}
{\fontsize{9.4}{12.6}\selectfont\color{fg2}Before anything expensive, try something cheap. Flatten the sequence and run it through two dense layers.\par}\vskip 4pt
\icode{python}{listing 13.10 --- a densely connected model}{listings/ch13-07.py}
\begin{iband}[signal]
\textbf{MSE as the loss, MAE as the metric.} MSE is used for training because it is \hilite{smooth around zero}, which gradient descent needs; MAE is reported because it is the number a human can interpret.
\end{iband}
\end{frame}

\framekicker{Section 13.2.3}
\begin{frame}[shrink=25]{\ldots{}and it does not beat the baseline}
\begin{iband}[rose]
Some validation losses come close to the no-learning baseline, \textbf{but not reliably}. This is exactly why the baseline was worth computing: it turns out to be \hilite{not easy to outperform}
\end{iband}
{\fontsize{9.4}{12.6}\selectfont\color{fg2}As the book puts it: \textbf{your common sense contains a lot of valuable information to which a machine learning model has no access.}\par}\vskip 4pt
\end{frame}

\framekicker{Section 13.2.3}
\begin{frame}[shrink=25]{Why gradient descent could not find the obvious answer}
{\fontsize{9.4}{12.6}\selectfont\color{fg2}A natural objection: the common-sense heuristic \textit{is} expressible as a two-layer network. So why does training not find it?\par}\vskip 4pt
\begin{iband}[signal]
Because your \textbf{hypothesis space} is the space of all two-layer networks with that configuration, and the heuristic is \textbf{one model among millions} in it. \hilite{Like looking for a needle in a haystack.}
\end{iband}
{\fontsize{9.4}{12.6}\selectfont\color{fg2}The general limitation: \textbf{unless the learning algorithm is hardcoded to look for a specific kind of simple model, it can fail to find a simple solution to a simple problem.}\par}\vskip 4pt
\end{frame}
\note{This is one of the most quotable passages in the book. It is the clearest argument for feature engineering and architecture priors anywhere in it.}

\framekicker{Section 13.2.4}
\begin{frame}[shrink=25]{Attempt 2 --- a 1D ConvNet}
{\fontsize{9.4}{12.6}\selectfont\color{fg2}The input has daily cycles, so a convolutional prior looks reasonable: the same pattern recurring at different points in time.\par}\vskip 4pt
\icode{python}{the 1D convolutional model}{listings/ch13-08.py}
\ioutput{listings/ch13-09.txt}
\end{frame}

\framekicker{Section 13.2.4}
\begin{frame}[shrink=25]{Worse than the dense model. Two reasons.}
\begin{minipage}[t]{0.494\textwidth}
\begin{icardbad}
\cardh{Weather is not translation invariant}
There are daily cycles, but \textbf{morning data behaves differently from evening or night data}. The invariance holds only at one very specific timescale.
\end{icardbad}
\end{minipage}%
\hfill
\begin{minipage}[t]{0.494\textwidth}
\begin{icardbad}
\cardh{Order matters, a lot}
The recent past is far more informative than five days ago --- and \textbf{max pooling and global average pooling destroy order information}.
\end{icardbad}
\end{minipage}%
\begin{iband}[signal]
A worked demonstration of chapter 9's claim: an architecture encodes assumptions, and \hilite{when those assumptions are false the architecture actively hurts}
\end{iband}
\end{frame}

\coursesection{02}{Recurrent neural networks}{Treat the data as what it is: a sequence where order matters.}

\framekicker{Section 13.3.1}
\begin{frame}[shrink=25]{Every network so far had no memory}
\begin{center}
\includegraphics[width=0.92\linewidth,height=0.60\textheight,keepaspectratio]{figs/ch13-feedvsrec.tex.pdf}
\end{center}
\vskip 3pt{\color{ink3}\fontsize{7.4}{9.4}\selectfont The dense model removed time; the convolutional model destroyed order. Neither treated the data as a sequence.\par}
\end{frame}

\framekicker{Section 13.3.1}
\begin{frame}[shrink=25]{The biological analogy the book uses}
{\fontsize{9.4}{12.6}\selectfont\color{fg2}As you read this sentence you process it \textbf{word by word --- or rather, eye saccade by eye saccade --- while keeping memories of what came before}. That gives you a fluid representation of its meaning.\par}\vskip 4pt
\begin{iband}[signal]
Biological intelligence processes information \textbf{incrementally, while maintaining an internal model} built from past information and constantly updated. An RNN adopts the same principle --- \hilite{in an extremely simplified version}
\end{iband}
\end{frame}

\framekicker{Section 13.3.1}
\begin{frame}[shrink=25]{An RNN is a network with an internal loop}
\begin{center}
\includegraphics[width=0.92\linewidth,height=0.60\textheight,keepaspectratio]{figs/ch13-rnnloop.tex.pdf}
\end{center}
\vskip 3pt{\color{ink3}\fontsize{7.4}{9.4}\selectfont Figure 13.6 --- it iterates through the sequence, keeping a state describing what it has seen so far.\par}
{\fontsize{9.4}{12.6}\selectfont\color{fg2}The state is \textbf{reset between independent sequences}, so one sequence is still one data point. What changes is that \hilite{the data point is no longer processed in a single step}\par}\vskip 4pt
\end{frame}

\framekicker{Section 13.3.2}
\begin{frame}[shrink=25]{Why you will rarely use SimpleRNN}
{\fontsize{9.4}{12.6}\selectfont\color{fg2}Keras has \inlinecode{SimpleRNN}, and it is \textbf{generally too simplistic to be of real use}. In theory it can retain information from many timesteps back; in practice such long-term dependencies prove impossible to learn.\par}\vskip 4pt
\begin{iband}[rose]
The cause is \textbf{the vanishing gradient problem} --- the same effect chapter 9 described for very deep feedforward networks. Studied by Hochreiter, Schmidhuber and Bengio in the early 1990s.
\end{iband}
{\fontsize{9.4}{12.6}\selectfont\color{fg2}Keras offers two alternatives designed to fix it: \textbf{\inlinecode{LSTM}} and \textbf{\inlinecode{GRU}}.\par}\vskip 4pt
\end{frame}

\framekicker{Section 13.3.2}
\begin{frame}[shrink=25]{LSTM: a conveyor belt running beside the sequence}
{\fontsize{9.4}{12.6}\selectfont\color{fg2}Long Short-Term Memory, from Hochreiter and Schmidhuber (1997), was the culmination of their work on vanishing gradients.\par}\vskip 4pt
\begin{iband}[signal]
Imagine a \textbf{conveyor belt running parallel to the sequence}. Information can jump onto it at any point, be carried to a later timestep, and jump off \textbf{intact} when needed --- \hilite{preventing older signals from gradually vanishing}
\end{iband}
{\fontsize{9.4}{12.6}\selectfont\color{fg2}The book points out this should remind you of \textbf{residual connections} from chapter 9. It is very nearly the same idea, applied along time rather than along depth.\par}\vskip 4pt
\end{frame}

\framekicker{Section 13.3.2}
\begin{frame}[shrink=25]{The carry track, in a picture}
\begin{center}
\includegraphics[width=0.92\linewidth,height=0.60\textheight,keepaspectratio]{figs/ch13-lstm.tex.pdf}
\end{center}
\vskip 3pt{\color{ink3}\fontsize{7.4}{9.4}\selectfont Figures 13.8--13.9 --- the carry \inlinecode{c} modulates both the next output and the next state.\par}
{\fontsize{9.4}{12.6}\selectfont\color{fg2}That is all the structural difference amounts to: a \inlinecode{SimpleRNN} cell, plus \textbf{an extra data flow that crosses timesteps without being squashed at every step}.\par}\vskip 4pt
\end{frame}

\framekicker{Section 13.3.2}
\begin{frame}[shrink=25]{And it finally beats the baseline}
{\fontsize{9.4}{12.6}\selectfont\color{fg2}One layer, sixteen units. Nothing about it is elaborate --- the architecture prior is what changed.\par}\vskip 4pt
\icode{python}{a simple LSTM model}{listings/ch13-10.py}
\begin{minipage}[t]{0.494\textwidth}
\istat{2.44}{common-sense baseline, validation MAE}
\end{minipage}%
\hfill
\begin{minipage}[t]{0.494\textwidth}
\istat{2.36}{LSTM, validation MAE}
\end{minipage}%
\end{frame}

\framekicker{Sections 13.2 -- 13.3}
\begin{frame}[shrink=25]{The whole chapter, as one ladder}
\begin{center}
\includegraphics[width=0.92\linewidth,height=0.60\textheight,keepaspectratio]{figs/ch13-ladder.tex.pdf}
\end{center}
\vskip 3pt{\color{ink3}\fontsize{7.4}{9.4}\selectfont Two deep learning models lose to the heuristic before one beats it.\par}
\begin{iband}[signal]
The margin over the baseline is \textbf{small}, and the chapter does not pretend otherwise. The lesson is not that LSTMs are powerful; it is that \hilite{the right architecture prior is what made the difference}
\end{iband}
\end{frame}

\coursesection{03}{Getting more out of RNNs}{Recurrent dropout, stacking, and bidirectional layers.}

\framekicker{Section 13.3.4}
\begin{frame}[shrink=25]{Dropout in a recurrent layer is not obvious}
{\fontsize{9.4}{12.6}\selectfont\color{fg2}The LSTM overfits quickly despite having very few units: training and validation losses diverge after a few epochs. Dropout is the obvious remedy --- \textbf{but applying it naively makes things worse}.\par}\vskip 4pt
\begin{iband}[amber]
It has long been known that applying dropout \textbf{before} a recurrent layer \hilite{hinders learning rather than helping regularisation}
\end{iband}
\end{frame}

\framekicker{Section 13.3.4}
\begin{frame}[shrink=25]{Yarin Gal's result: hold the mask constant}
\begin{center}
\includegraphics[width=0.92\linewidth,height=0.60\textheight,keepaspectratio]{figs/ch13-dropoutmask.tex.pdf}
\end{center}
\vskip 3pt{\color{ink3}\fontsize{7.4}{9.4}\selectfont The same pattern of dropped units at every timestep, rather than a fresh random one.\par}
{\fontsize{9.4}{12.6}\selectfont\color{fg2}Determined in 2015 as part of Gal's PhD thesis on Bayesian deep learning. He did the research \textbf{using Keras}, and helped build the mechanism directly into its recurrent layers.\par}\vskip 4pt
\end{frame}

\framekicker{Section 13.3.4}
\begin{frame}[shrink=25]{The two arguments that implement it}
{\fontsize{9.4}{12.6}\selectfont\color{fg2}Every recurrent layer in Keras takes two dropout-related arguments, and they regularise different things.\par}\vskip 4pt
\icode{python}{recurrent dropout}{listings/ch13-11.py}
\begin{itable}{>{\raggedright\arraybackslash}p{0.282\textwidth}>{\raggedright\arraybackslash}p{0.658\textwidth}}
\thead{Argument} & \thead{What it drops} \\ \midrule
\inlinecode{dropout} & The layer's \textbf{input} units. \\
\inlinecode{recurrent\_dropout} & The layer's \textbf{inner recurrent activations} --- a temporally constant mask, which is the part Gal's result is about. \\
\end{itable}
\end{frame}

\framekicker{Section 13.3.5}
\begin{frame}[shrink=25]{Stacking: when you stop overfitting, add capacity}
{\fontsize{9.4}{12.6}\selectfont\color{fg2}Once dropout has stopped the overfitting but performance has plateaued, chapter 5's rule applies: \textbf{increase capacity until overfitting becomes the main obstacle again}.\par}\vskip 4pt
\begin{iband}[signal]
Recurrent layer stacking is the classic way to do it. Not long ago \textbf{Google Translate was powered by a stack of seven large LSTM layers} --- \hilite{which is huge}
\end{iband}
\end{frame}

\framekicker{Section 13.3.5}
\begin{frame}[shrink=25]{The one argument that makes stacking work}
\begin{center}
\includegraphics[width=0.92\linewidth,height=0.60\textheight,keepaspectratio]{figs/ch13-stack.tex.pdf}
\end{center}
\vskip 3pt{\color{ink3}\fontsize{7.4}{9.4}\selectfont Intermediate layers must emit a full sequence; only the last one emits a single vector.\par}
\icode{python}{a stack of two GRU layers}{listings/ch13-12.py}
\begin{iband}[signal]
\inlinecode{return\_sequences=True} makes a layer return a \textbf{rank-3 tensor} --- its output at every timestep --- instead of only the last one. Omit it on an intermediate layer and \hilite{the next layer has no sequence to read}
\end{iband}
\end{frame}

\framekicker{Section 13.3.5}
\begin{frame}[shrink=25]{GRU, and why the book switches to it here}
{\fontsize{9.4}{12.6}\selectfont\color{fg2}A \textbf{Gated Recurrent Unit} is a streamlined variant of LSTM: the same idea of a carried track, with a simpler gating structure and therefore fewer parameters.\par}\vskip 4pt
\begin{iband}[signal]
In practice the choice between them is empirical. GRU is \textbf{cheaper to run}; LSTM has \textbf{slightly more representational power}. \hilite{Try both and measure} rather than arguing about it.
\end{iband}
\end{frame}

\framekicker{Section 13.3.6}
\begin{frame}[shrink=25]{Bidirectional layers: the same data, read backwards too}
{\fontsize{9.4}{12.6}\selectfont\color{fg2}A bidirectional RNN runs \textbf{two copies} of a recurrent layer --- one over the sequence in order, one in reverse --- and merges their representations.\par}\vskip 4pt
\icode{python}{wrapping a layer}{listings/ch13-13.py}
\begin{iband}[amber]
On \textbf{this} problem it does not help: for weather, the recent past is what matters, so reading in reverse gives a \hilite{chronologically backwards view that is simply worse}. It shines on text, where both directions carry meaning --- chapter 14 returns to it.
\end{iband}
\end{frame}

\framekicker{Section 13.4}
\begin{frame}[shrink=25]{Going further}
\begin{ibullets}
\item \textbf{Adjust the number of units} in each recurrent layer, and the dropout rates --- the configurations here were chosen largely arbitrarily.
\item \textbf{Adjust the learning rate} of the optimizer.
\item Try a \textbf{stack of \inlinecode{Dense} layers} as the regressor on top instead of a single one.
\item \textbf{Improve the input}: try longer or shorter sequences, or a different sampling rate.
\end{ibullets}
\begin{iband}[signal]
And the standing warning: \textbf{always evaluate on the validation set}, then confirm once on test --- otherwise you end up \hilite{overfitting to your validation procedure}, exactly as chapter 5 described.
\end{iband}
\end{frame}

\framekicker{Practice}
\begin{frame}[shrink=25]{What this chapter changes about how you start}
\begin{isteps}
\item \textbf{Compute the baseline before writing any model.} It is cheap, and here it would have saved two failed attempts.
\item \textbf{State the architecture prior out loud.} \textit{Convolution assumes translation invariance} --- is that true of your data, at the timescale you care about?
\item \textbf{Check whether order matters.} If it does, any layer that pools or flattens is throwing away the signal.
\end{isteps}
\begin{iband}[signal]
For transactional or operational data these three questions usually settle the architecture before a single experiment runs --- which is \hilite{the cheapest hour in the whole project}
\end{iband}
\end{frame}

\framekicker{Summary}
\begin{frame}[shrink=25]{What has to survive this chapter}
\begin{isteps}
\item \textbf{Always establish a common-sense baseline first.} Here it beat two deep learning models.
\item \textbf{A good solution existing in your hypothesis space does not mean gradient descent will find it.}
\item \textbf{Dense flattens away time; 1D convolution pools away order.} Neither prior fits a forecasting problem.
\item \textbf{RNNs keep a state} and process a sequence incrementally.
\item \textbf{\inlinecode{SimpleRNN} cannot learn long dependencies}; LSTM adds a carry track --- residual connections, along time.
\end{isteps}
\end{frame}

\framekicker{Summary}
\begin{frame}[shrink=25]{\ldots{}and the three refinements}
\begin{minipage}[t]{0.3253\textwidth}
\begin{icardaccent}
\cardh{Recurrent dropout}
A \textbf{temporally constant mask}. A fresh random mask each timestep disrupts the error signal through time.
\end{icardaccent}
\end{minipage}%
\hfill
\begin{minipage}[t]{0.3253\textwidth}
\begin{icardaccent}
\cardh{Stacking}
More capacity, at more compute. Intermediate layers need \inlinecode{return\_sequences=True}.
\end{icardaccent}
\end{minipage}%
\hfill
\begin{minipage}[t]{0.3253\textwidth}
\begin{icardaccent}
\cardh{Bidirectional}
Helps where \textbf{both directions carry meaning} --- text. Does not help where only the recent past matters.
\end{icardaccent}
\end{minipage}%
\vskip 4pt
\reslink{NOTEBOOK}{03\_lstm\_and\_gru.ipynb}{https://hendrikarisma.my.id/rs/ai-products-course/notebooks/ch13/03\_lstm\_and\_gru.html}
\reslink{NEXT}{Chapter 14 --- Text classification}{https://hendrikarisma.my.id/rs/ai-products-course/ch14/index.html}
\vskip 2pt
\end{frame}

\end{document}
